% Avsnitt 19.3-19.8, ur lectures/L19/appendix/b_fpga_bringup_and_review.md och lectures/L19/README.md.

\section{Att konstruera \code{can_spi_node}}
\label{c19:sec:node}

\code{can_controller} är \textbf{inte} den entitet som ska upp på kortet, och det är värt att vara
precis med. Dess gränssnitt ligger på registernivå: \code{tx_id} är 11 bitar, \code{tx_dlc} 4,
\code{tx_data} 64, plus \code{tx_req} och fem statusutgångar. Ingenting på ett DE0-CV erbjuder ett så
brett gränssnitt, och i det färdiga systemet kommer de ingångarna från drivrutinens register snarare
än från en mänsklig hand.

\Chapref{c18:ch} byggde lagret som gör dem nåbara, och \secref{c19:sec:bridge} transporten framför
det. Den sista modul boken skriver är alltså \file{can_spi_node.vhd}: toppnivån som binder ihop
\code{spi_slave}, \code{spi_reg_bridge}, \code{register_bank} och \code{can_controller} till en nod,
och exponerar exakt två saker mot omvärlden, SPI och CAN-bussen.

Till skillnad från varje annan modul i projektet kommer den utan egen testbänk på modulnivå, och det
är avsiktligt: varje block inuti den har redan en, och det som återstår \textemdash{} pinnar,
kopplingar och spänningsnivåer \textemdash{} är precis det en simulering inte kan modellera. Kortet
är dess testbänk, och bring-up-stegen i \secref{c19:sec:bringup} är hur man läser den. (Det finns en
utdelad systemtestbänk, \code{can_spi_node_tb}, som driver två kompletta noder från modellerade
SPI-mastrar; den är den enda i boken som driver \code{sclk}, \code{mosi} och \code{ss} som riktiga
vågformer.)

\subsection*{Gränssnitt}

\begin{booktable}{@{}p{4.6em}p{3.4em}p{5.2em}L@{}}
  \thead{Port} & \thead{Riktn.} & \thead{Typ} & \thead{Betydelse} \\
  \midrule
  \code{clock} & in & \code{std_logic} & 50 MHz-oscillatorn på kortet. Konstruktionens enda klocka. \\
  \code{reset_n} & in & \code{std_logic} & Rå, aktiv låg reset, direkt från en pinne. Asynkron; den
    här nivån synkroniserar den. \\
  \code{sclk} & in & \code{std_logic} & SPI-klockan från mastern. Används aldrig som klocka. \\
  \code{mosi} & in & \code{std_logic} & Ut från mastern, in i slaven. \\
  \code{ss} & in & \code{std_logic} & Chip select, aktiv låg. Ramar in transaktionen. \\
  \code{miso} & out & \code{std_logic} & Ut från slaven, in i mastern. \\
  \code{rx_bus} & in & \code{std_logic} & CAN-bussens nivå, direkt från pinnen och asynkron.
    \code{can_controller} synkroniserar den själv. \\
  \code{tx_bus} & out & \code{std_logic} & Biten den här noden vill driva. \\
  \code{bus_en} & out & \code{std_logic} & Om den driver över huvud taget. \\
\end{booktable}

De tre bussportarna är det gränssnitt på digitalsidan som en riktig CAN-transceiver presenterar
(\chapref{c10:ch}), och det är därför de lämnar konstruktionen som tre separata signaler i stället
för en ledning. Att göra en enda fysisk ledning av dem är ett jobb på kortnivå, och
\secref{c19:sec:wire} handlar om det.

Ordningen är SPI:s fyra ledningar som en grupp, \code{sclk}/\code{mosi}/\code{ss}/\code{miso}, och
därefter CAN-bussens tre. \code{miso} ligger alltså bland SPI-portarna och inte hos de andra
utgångarna, så att de fyra pinnarna står i samma ordning som i kopplingstabellen i
\secref{c19:sec:bringup}. Det spelar roll: \code{can_spi_node_tb} binder positionellt som allt annat,
och \code{miso} och \code{rx_bus} har samma typ, så en förväxling av dem passerar analysen utan ett
ord och ger en nod som aldrig svarar.

Det finns inga generics. Allt konstruktionen behöver som en pinne inte kan leverera är antingen en
konstant i \code{can_def} eller ett register som drivrutinen skriver.

\subsection*{Beteende: strukturellt, och det är poängen}
\code{can_spi_node} har ingen tillståndsmaskin, ingen protokollkunskap och nästan ingen logik alls.
Den har exakt tre uppgifter:
\begin{itemize}
  \item \textbf{Instansiera de fyra blocken och koppla dem på rad.} \code{spi_slave} gör om
    SCK/MOSI/SS till bytes; \code{spi_reg_bridge} gör om bytes till registeråtkomster;
    \code{register_bank} gör om registeråtkomster till signaler på kontrollernivå;
    \code{can_controller} gör om dem till en ram. Varje gräns är ett av de kontrakt som redan
    verifierats av en utdelad testbänk, så om kedjan beter sig illa sitter felet i en koppling här,
    inte i ett block.
  \item \textbf{Synkronisera reseten.} En enda \code{meta_prev}-instans tar den råa \code{reset_n}
    och producerar den \code{reset_s2_n} som varje delblock tar. \code{can_controller} är undantaget
    som bekräftar regeln: den äger sin resetsynkroniserare internt, vid sidan av den för
    \code{rx_bus}, så den här nivån matar den med den råa \code{reset_n} och synkroniserar bara åt
    blocken på SPI-sidan. Det ger tre \code{meta_prev}-instanser i hela konstruktionen \textemdash{}
    två inuti \code{can_controller}, en här.
  \item \textbf{Låt portarna mot kontrollern gå rakt över.} \code{register_bank}s \code{tx_id},
    \code{tx_dlc}, \code{tx_data} och \code{tx_req} till \code{can_controller}, och dess
    \code{tx_done}, \code{rx_id}, \code{rx_dlc}, \code{rx_data}, \code{rx_valid} och \code{error}
    tillbaka. Port till port, ingenting emellan. \Chapref{c18:ch} konstruerade de två gränssnitten så
    att de speglar varandra just för att det ska stämma.
\end{itemize}

Om en regel i den här filen någonsin behöver nämna ett CAN-fält eller en SPI-byte hör den hemma i en
annan fil. Det är testet att tillämpa medan du skriver den.

% ------------------------------------------------------------------------------------------
\section{Att få ut bussen på en ledning}
\label{c19:sec:wire}

\code{tx_bus} och \code{bus_en} är ett \emph{par}, inte en signal: konstruktionen säger ”driv
dominant”, ”släpp”, aldrig ”driv recessivt”. Två noders push-pull-utgångar knutna till en ledning
skulle slåss med varandra i samma ögonblick som den ena drev låg och den andra hög. Två sätt att få
den delade ledningen att uppföra sig:

\paragraph{Med en transceiver}
Vilket är vad en riktig nod använder: \code{tx_bus} går till dess TXD, dess RXD kommer tillbaka som
\code{rx_bus}, och transceivern sköter arbetet med öppen dränering analogt.

\paragraph{Utan en}
Med två DE0-CV kopplade direkt ihop för tvånodsdemot gör FPGA:ns egen I/O jobbet. Välj en GPIO-pinne
och driv den med öppen dränering:

\begin{vhdlcode}
GPIO_0(2) <= '0' when ((bus_en = '1') and (tx_bus = '0')) else 'Z';
rx_bus    <= GPIO_0(2);
\end{vhdlcode}

med en pull-up som återställer \code{'1'} medan varje nod släppt ledningen \textemdash{} FPGA:ns
svaga pull-up aktiverad på den pinnen i Pin Planner, eller ett externt motstånd. Det är
\secref{c10:sec:can}:s wired-AND-regel gjord fysisk, i en enda konkurrent tilldelning.
\code{GPIO_0(0)} och \code{GPIO_0(1)} är värda att driva med \code{tx_bus} och \code{bus_en}
oregistrerade också, enbart för att en logikanalysator ska kunna se vad noden \emph{avsåg} vid sidan
av vad ledningen faktiskt gjorde; varje oanvänt index drivs \code{'Z'}.

Ingenting här synkroniseras, av samma skäl som \code{reset_n} inte synkroniseras innan den når
\code{can_controller}: kontrollern äger sina egna asynkrona ingångar, så pinnen dras rakt till porten
\code{rx_bus}.

Pinnindexen är en bekvämlighet i Pin Planner, inte ett krav. Vilka tre pinnar som helst på listen
fungerar, så länge två ihopkopplade kort är överens; se bilaga~\ref{app:quartus}.

% ------------------------------------------------------------------------------------------
\section{Valfritt: en wrapper på kortnivå för felsökning}
\label{c19:sec:wrapper}

Allt ovanstående går bara att nå över SPI, vilket betyder att när länken inte fungerar har du ingen
möjlighet att se in i den. En liten wrapper med knappar och lysdioder, som står bredvid
\code{can_spi_node} i stället för att ersätta den, är värd en timme om bring-upen kör fast
\textemdash{} och de tre beslut den tvingar fram är värda att förstå oavsett om du bygger den eller
inte.

Wrappern driver \code{register_bank}s portar direkt från knappar och omkopplare och låser dess
statusutgångar på lysdioder:

\lecfig[0.86]{lectures/L19/appendix/images/can_demo.png}{En wrapper på kortnivå: knappar och
omkopplare in i registerbanken, statusbitarna ut på lysdioder.}{c19:fig:demo}

\begin{itemize}
  \item \textbf{En hållen knapp måste bli en encykelspuls}, och inte bara för snygghetens skull. En
    knapp som hålls i några millisekunder är hundratusentals klockcykler vid 50~MHz. Mata in det rakt
    i \code{tx_req} så startar kontrollern en ny ram i samma ögonblick som den förra tar slut, om och
    om igen \textemdash{} och varje passage genom starttillståndet suddar \code{error} innan något
    hunnit läsa det. Alltså: \textbf{synkronisera} den råa pinnen med en \code{meta_prev},
    \textbf{studsfiltrera} den synkroniserade nivån med en räknare som accepterar ett nytt värde
    först när det legat kvar i ungefär 10~ms, och \textbf{därefter} flankdetektera den. Alla tre, i
    den ordningen: att flankdetektera före studsfiltreringen ger en puls per studs, och att
    studsfiltrera före synkroniseringen matar räknaren med en signal som aldrig samplats säkert. Det
    här är samma konvertering som \code{register_bank} utför för \code{TX_SEND} (\chapref{c18:ch}),
    och det är därför banken gör wrappern valfri snarare än nödvändig.
  \item \textbf{En omkopplarbank är bedömningsfrågan.} Tio omkopplare som korsar in i klockdomänen
    synkroniseras bit för bit, vilket tar bort metastabiliteten men inte \emph{skevheten}: två
    omkopplare som slås om samtidigt kan landa en klocka isär, så \code{tx_id} skulle kortvarigt
    kunna läsa ett värde som aldrig stod på omkopplarna. Det är acceptabelt bara därför att en
    människa ställer in dem och sedan låter dem vara. Koppla samma tio bitar till något som ändras
    varje cykel och en bit-för-bit-synkroniserare är helt fel verktyg \textemdash{} då krävs en
    handskakning, eller en kodning där bara en bit ändras i taget. Värt att säga rakt ut, eftersom
    ”bredda bara synkroniseraren” är en vana som förr eller senare biter ifrån.
  \item \textbf{Statuslysdioder behöver låsas}, eftersom \code{tx_done} och \code{rx_valid} är pulser
    på 20~ns. Det är samma problem som \code{register_bank} löser åt drivrutinen, och att lösa det
    två gånger på två ställen är precis så de två kopiorna glider isär \textemdash{} så driv
    lysdioderna från bankens STATUS-bitar i stället för från kontrollerns pulser, så förblir wrappern
    en wrapper.
\end{itemize}

% ------------------------------------------------------------------------------------------
\section{Bring-up}
\label{c19:sec:bringup}

Bring-upen har två halvor: CAN-sidan mellan två kort, och SPI-sidan mellan en mikrokontrollerhatt och
kortet. Ta dem i den ordningen. CAN-halvan behöver ingen mikrokontroller alls, så den går att köra
medan hatten fortfarande ligger på lödbordet, och den utesluter en hel felklass innan någon
SPI-ledning finns att misstänka.

När konstruktionen passerar simuleringen fullt ut, syntetisera och programmera \code{can_spi_node}:
New Project Wizard, enheten \code{5CEBA4F23C7N}, Pin Planner, Compile, Programmer. Stegen står i
bilaga~\ref{app:quartus}, som är den auktoritativa referensen för det flödet, ända ner till tillägget
för Cyclone V-enheter och till att få USB-Blastern synlig för Programmer. Notera att \code{CLOCK_50}
är den enda pinne du inte får välja fritt.

\subsection*{CAN-halvan: två kort på en buss}
Två kort, deras busslinjer på \code{GPIO_0(2)} knutna ihop till \textbf{en} ledning med gemensam
jord, och drivningen med öppen dränering ur \secref{c19:sec:wire} gör den delade ledningen till ett
wired-AND.

\begin{enumerate}
  \item \textbf{Ett ensamt kort.} Busslinjen bär ingenting annat än pull-upen, så noden läser
    tillbaka exakt den nivå den driver: en loopback. \code{tx_bus} och \code{bus_en} ska växla fält
    för fält på en logikanalysator via \code{GPIO_0(0)}/\code{GPIO_0(1)}, \code{tx_done} pulsa vid
    slutet av EOF, och \code{error} aldrig gå hög.
  \item \textbf{Två kort, en ledning.} Tryck på en knapp på det ena kortet och en lysdiod ska tändas
    på det andra. Det är första gången konstruktionen möter två \emph{oberoende} oscillatorer; i
    \code{can_controller_tb} delade noderna en \code{clock}.
  \item \textbf{Arbitrering på riktigt.} Låt båda korten sända samtidigt och se att den med lägst
    identifierare vinner, och att förloraren backar.
\end{enumerate}

En loopback är också där \secref{c10:sec:bittiming}:s sampelpunkt på 70~\% i tysthet gör rätt för
sig. Nodens egen signal tar nu en riktig rundtur, ut genom en pad, längs ledningen, in igen, och genom
\code{rx_bus}-synkroniserarens två vippor, till en kostnad av en handfull klockcykler. Med
\code{TICKS_PER_BIT = 50} och \code{SAMPLE_TICK = 35} läser noden bussen 35 tick in i en bit den
började driva vid tick 0, så dess eget eko har sedan länge stabiliserat sig och
arbitreringsövervakningen läser aldrig fel på det. Att sampla vid bitperiodens \emph{början} skulle i
stället läsa ekot av biten före, och arbitreringsövervakningen (sänd recessivt, läs dominant,
\chapref{c17:ch}) skulle rapportera en falsk förlust vid den första recessiva bit som följer på en
dominant.

En loopback kan medvetet inte visa noden \emph{ta emot} den ramen. \code{can_controller} låser
\code{role = '1'} för hela ramen och \code{rx_shift_reg} går bara medan \code{role = '0'}
(\chapref{c17:ch}), så en sändare avkodar aldrig sin egen trafik; det är precis därför
systemtestbänken behöver två noder.

\textbf{Det andra kortet sluter den lucka \secref{c19:sub:gaps-tb} kallar den största.} Testbänkens
två noder delar en \code{clock}: \code{resync} fasriktar dem fortfarande vid varje SOF, men deras
timrar kan aldrig \emph{driva isär} mitt i ramen. Två kort har två kristaller, som går på 50~MHz bara
inom den tolerans komponenterna är specificerade till, så deras bitperioder skiljer sig verkligen åt,
och mottagaren frikör på sin egen timer resten av ramen efter sin enda \code{resync} vid SOF. Ärlig
aritmetik håller förväntningarna i schack: vid typiska kristalltoleranser på några tiotals ppm är
driften över en ram en liten bråkdel av en tick, så det demot synligt bevisar är att hela kedjan
fungerar över genuint oberoende klockor från en godtycklig startfas; hur mycket drift konstruktionen
skulle tåla innan sampelpunkten vandrar av biten är beräkningen i Övning~\ref{c12:ex:drift}.

\subsection*{SPI-halvan: en mikrokontrollerhatt mot kortet}
Kopplingen, per nod. Den är inte din att välja: den står i bilaga~\ref{app:protocol}, eftersom
parallellklassen skriver sin drivrutin mot samma tabell.

\begin{booktable}{@{}p{6em}p{6em}L@{}}
  \thead{Signal} & \thead{AVR32DB28} & \thead{DE0-CV} \\
  \midrule
  \code{SCK} & \code{PC2} & \code{sclk} \\
  \code{MOSI} & \code{PC0} & \code{mosi} \\
  \code{MISO} & \code{PC1} & \code{miso} \\
  \code{SS} & \code{PC3} & \code{ss} \\
  I/O-matning & \code{VDDIO2} & 3,3 V \\
  Kärnmatning & \code{VDD} & (hattens egen 5 V) \\
  GND & \code{GND} & GND \\
\end{booktable}

Ingen nivåomvandlare behövs, och det är ett val och inte en förenkling: AVR32DB28:ans \code{PORTC} är
en egen spänningsdomän, matad från \code{VDDIO2} i stället för \code{VDD}, så den läggs på DE0-CV:ns
3,3~V medan kärnan går kvar på 5~V. Varje nivå FPGA:n ser är då en nivå den är specificerad för, och
de 3,3~V FPGA:n driver tillbaka läses mot en 3,3 V-tröskel i stället för en 5 V-tröskel. SPI0 flyttas
dit med \code{PORTMUX.SPIROUTEA = ALT1}.

\begin{warning}
\textbf{Tre krav på hatten, och de måste vara på plats innan den etsas.}
\begin{itemize}
  \item \textbf{\code{PC0}--\code{PC3} är reserverade.} De är SPI0:s ALT1-mux och får inte bära någon
    annan kringutrustning. Ingen lysdiod, ingen knapp, och framför allt ingen avkopplingskondensator
    på \code{SS} \textemdash{} en knapp med 100~nF på pinnen gör \code{SS}-flanken obrukbar vid
    1~MHz.
  \item \textbf{Ingenting annat på \code{PORTC}.} Hela porten ligger i \code{VDDIO2}-domänen på
    3,3~V, så en 5 V-kringutrustning där skulle både sluta fungera och mata 5~V in i en 3,3 V-domän.
  \item \textbf{\code{VDDIO2} är en egen bana.} Den får inte knytas till \code{VDD} på kortet, utan
    ska ut till en pinne där DE0-CV:ns 3,3~V kopplas in, med 100~nF nära chippet.
\end{itemize}

Att domänen över huvud taget är separat styrs dessutom av \code{MVSYSCFG} i \code{FUSE.SYSCFG1}: i
enkelmatningsläge knyts \code{VDDIO2} internt till \code{VDD}, \code{PORTC} blir en 5 V-port, och
kopplingen ovan lägger då 5~V rakt in i ingångar som inte är 5 V-toleranta \textemdash{} på en hatt
som ser precis likadan ut som en rätt byggd. Kontrollera fusen före första inkopplingen. Det är det
enda felet i det här kapitlet som kan kosta ett kort.
\end{warning}

Stegen tas i den här ordningen, och inte i någon annan. Varje steg utesluter en felkälla, så att det
som fortfarande inte fungerar efter steg $n$ med säkerhet ligger i steg $n+1$:

\begin{enumerate}
  \setcounter{enumi}{3}
  \item \textbf{SPI-loopback på enbart hatten}, med \code{PC1} kopplad till \code{PC0} och FPGA:n
    urkopplad. Bevisar att mastern fungerar, och på köpet att \code{PORTMUX} och \code{VDDIO2} är
    rätt satta.
  \item \textbf{Registereko}: skriv \code{TX_ID}, läs tillbaka det. Bevisar att hela
    transaktionskedjan bär.
  \item \textbf{\code{STATUS}-pollning}: läs \code{STATUS} och se att bit 0 är satt efter reset.
  \item \textbf{En nod sänder på kommando}: skriv en ram och utlös \code{TX_SEND}, och se
    \code{STATUS} bit 0 gå låg och tillbaka \textemdash{} nu utlöst över SPI i stället för av en
    knapp.
  \item \textbf{Två noder, hela kedjan}: den ena sänder på SPI-kommando, den andra tar emot och
    kvitterar med \code{RX_ACK}. Det är systemet.
\end{enumerate}

\paragraph{Vad som återstår}
När passet är slut har båda halvorna gått på riktig hårdvara. Det som återstår är inte kopplingen
utan \textbf{motparten}: din nod har svarat på ett testprogram som är skrivet för att lyckas, inte på
en drivrutin skriven av någon annan mot enbart specifikationen. Det är därför
protokollspecifikationens tidskrav \textemdash{} de 60~ns kring \code{SS}, regeln om vad \code{MISO}
bär under kommandobyten, och avbrottsregeln \textemdash{} är skrivna som krav och inte som råd: ett
testprogram som aldrig gör något oväntat bekräftar dem aldrig, och en drivrutin som gör det fäller dig
på dem.

% ------------------------------------------------------------------------------------------
\section{Hur produktionsklara CAN-kontrollrar går längre}
\label{c19:sec:gaps}

Den här bokens kontroller är medvetet förenklad: en undervisningskonstruktion snarare än en
certifierbar. Listan nedan är den här bokens version av den; \canref{7} ställer samma lista från
protokollhållet, tillsammans med blocköversikten och de tre begränsningar som når ända upp i
drivrutinens kod, och \canref{6} beskriver den felhantering \textemdash{} felramar, felräknare,
error-passive och bus-off \textemdash{} som den här konstruktionen helt saknar.

Anmärkningsvärda luckor, värda att namnge uttryckligen:
\begin{itemize}
  \item \textbf{Bara standardramar med data:} ledningsformatet är en bit-exakt standard-CAN-dataram,
    RTR inräknad, men det finns inga 29-bitars utökade identifierare (mottagaren kontrollerar inte
    ens IDE, så en utökad ram skulle tolkas fel tills dess CRC-kontroll misslyckas), och en fjärram
    (en recessiv RTR) avvisas med \code{error} i stället för att besvaras med de begärda data.
  \item \textbf{Inga felramar och ingen återhämtning från bus-off:} en riktig CAN-nod räknar
    sändnings- och mottagningsfel och går in i alltmer begränsade tillstånd (error-passive, bus-off)
    för att skydda en frisk buss från en trasig nod; den här bokens \code{can_controller} avbryter
    bara den pågående ramen lokalt.
  \item \textbf{Ingen detektering av bit- eller formfel:} den här sändaren jämför vad den sände mot
    vad bussen läste bara under arbitreringen (\chapref{c17:ch}). En riktig kontroller övervakar
    varje bit den sänder, flaggar för ett bitfel i samma ögonblick som bussen säger emot utanför
    arbitreringsfältet och ACK-luckan, och kontrollerar vid mottagning att svansens fasta fält har
    sina föreskrivna värden.
  \item \textbf{Ingen omsändning vid uteblivet kvitto:} en riktig kontroller försöker om automatiskt.
    Den här konstruktionen märker inte ens saken, eftersom sändaren aldrig läser tillbaka ACK-luckan.
  \item \textbf{Ingen acceptansfiltrering och ingen mottagningsbuffring:} varje ram på bussen landar
    i samma enda uppsättning \code{rx_*}-portar, och en andra ram som anländer innan drivrutinen läst
    den första skriver helt enkelt över den. Produktionsklara kontrollrar erbjuder identifierarfilter
    i hårdvara och en mottagnings-FIFO eller en uppsättning brevlådor. Det här är den lucka
    parallellklassens drivrutin känner mest direkt, och \chapref{c18:ch} visar exakt var den blir
    synlig på registernivå.
  \item \textbf{En fast modell för bittajmingen:} den här boken samplar vid en fast punkt på 70~\%;
    produktionsklara kontrollrar exponerar konfigurerbara propagerings- och fassegment och en
    synchronization jump width, avstämda efter busslängd och antal noder.
  \item \textbf{En avbruten mottagning (efter förlorad arbitrering) kastas helt enkelt}, i stället för
    att sömlöst fortsätta ta emot den ram som vann; en riktig nod fortsätter ta emot oavsett vem som
    vann bussen. Den här konstruktionens nod väntar ut resten av ramen i \code{STATE_IDLE}
    (\secref{c16:sec:roles}) och tar vid från nästa.
  \item \textbf{En förenklad regel om ledig buss:} den här konstruktionen kräver elva recessiva
    bitperioder, samma antal som riktig CAN begär av en nod som ansluter sig till busstrafiken, innan
    någon nod får lämna \code{STATE_IDLE}. Riktig CAN är mer finkornig: en nod som just deltagit i en
    ram får sända igen redan efter den 3-bitars intermission som följer på EOF, och en nod som varit
    bus-off väntar in 128 förekomster av elva recessiva bitar innan den får ansluta sig igen.
\end{itemize}

% ------------------------------------------------------------------------------------------
\section{Konstruktionen, sedd i backspegeln}
\label{c19:sec:review-design}

Varje kapitel från \chapref{c10:ch} och framåt lade till exakt ett ansvar, aldrig duplicerat någon
annanstans:

\begin{booktable}{@{}p{9.6em}Lp{5.4em}@{}}
  \thead{Block} & \thead{Äger} & \thead{Byggt i} \\
  \midrule
  \code{can_def} & Delade konstanter och subtyper som varje annat block är överens om &
    \chapref{c10:ch} \\
  \code{meta_prev} & Att få en asynkron signal säkert in i den här klockdomänen, utan någon aning om
    vad signalen betyder & \chapref{c12:ch} \\
  \code{bit_timer} & \emph{När} det ska samplas eller skiftas, utan någon aning om vad som sänds &
    \chapref{c12:ch} \\
  \code{crc15} & Beräkningen av CRC-15, en bit i taget, utan någon aning om vad en ”ram” är &
    \chapref{c13:ch} \\
  \code{tx_shift_reg} / \code{rx_shift_reg} & Bitstoppning och avstoppning, utan någon aning om vad
    ett ID eller en databyte är & \chapref{c14:ch}--\ref{c15:ch} \\
  \code{can_controller} & Det enda block som känner till det faktiska CAN-ramformatet, arbitreringen
    inräknad & \chapref{c11:ch} (portar), \chapref{c16:ch}--\ref{c17:ch} (beteende) \\
  \code{register_bank} & Pulser till pollbara nivåer, och registerkartans semantik; ingen aning om
    vad SPI eller CAN är & \chapref{c18:ch} \\
  \code{spi_slave} & Bytes av en seriell ledning, säkert in i 50 MHz-domänen; ingen aning om vad ett
    register är & utdelad, läses i \chapref{c18:ch} \\
  \code{spi_reg_bridge} & Fembytestransaktionen; ingen aning om vad något enskilt register betyder &
    \chapref{c19:ch} \\
  \code{can_spi_node} & Kopplingarna, och pinnarna bussen bor på. Ingen egen logik &
    \chapref{c19:ch} \\
\end{booktable}

Det är hela noden: fyra små block med ett ansvar var och ett delat paket komponerade inuti
\code{can_controller}, och sedan tre lager till framför den, där vart och ett bara känner till det
under sig. \code{can_spi_node} står utanför dem alla, eftersom att anpassa sig till ett visst kort är
en annan sak än att tala CAN, eller SPI, eller register.

Läs tabellen uppifrån och ned så är den också en lista över vad som skulle behöva ändras för att nå
ett annat system. Byt ut \code{spi_slave} och \code{spi_reg_bridge} mot ett I2C-par så är allt under
dem orört. Byt ut \code{can_controller} mot en annan protokollmotor så flyttar registerkartan men inte
transporten. Det är vad gränserna köpte.

\paragraph{Registerkartan i sammanfattning}
Varje port på registersidan av \code{can_controller} motsvarar ett register som parallellklassens
C++-drivrutin läser eller skriver; de återstående fem (\code{clock}, \code{reset_n} och de tre
bussignalerna) är fysiska och ingen drivrutin ser dem någonsin. Motsvarigheten är inte längre en plan,
som den var när \secref{c11:sec:regmap} först ritade den här tabellen \textemdash{}
\code{register_bank} är modulen som förverkligar den, så läs högerkolumnen som ”registret banken
presenterar vid det indexet”:

\begin{booktable}{@{}p{14em}L@{}}
  \thead{\code{can_controller}-port} & \thead{Register i kartan} \\
  \midrule
  \code{tx_id}, \code{tx_dlc}, \code{tx_data} & \code{TX_ID}, \code{TX_DLC},
    \code{TX_DATA_LO}/\code{HI} \\
  \code{tx_req}, \code{tx_done} & \code{TX_SEND}, en bit i \code{STATUS} \\
  \code{rx_id}, \code{rx_dlc}, \code{rx_data} & \code{RX_ID}, \code{RX_DLC},
    \code{RX_DATA_LO}/\code{HI} \\
  \code{rx_valid} & en bit i \code{STATUS}, nollställd via \code{RX_ACK} \\
  \code{error} & \code{ERROR_FLAGS} \\
\end{booktable}

Den motsvarigheten är där drivrutinsarbetet börjar: allt parallellklassen skriver mot står i
bilaga~\ref{app:protocol}, och ingenting i den här konstruktionens VHDL.
